\section{Tao and Logos}

A conference on the book \emph{Christ the Eternal Tao} by the author \textbf{Hieromonk Damascene} is available at Ancient Faith Radio\footnote{\url{http://ancientfaith.com/specials/christ\_the\_eternal\_tao}}.

\begin{quotex}
Learn how Eastern Orthodox Christian spirituality provides seekers of our day clear guidance on acquiring stillness, overcoming the passions, dealing with thoughts, and cultivating the virtues, as well as precise teachings on spiritual *deception, all of which* guides seekers more safely and surely on the path to communion with God (Tao --- Chinese). In the profound mystical and contemplative tradition of the Christian East, seekers are able to go well beyond the realizations in Eastern religions. A long--awaited answer for those who, having turned away from modern Western religiosity, are drawn to the freshness, directness and simplicity of Lao Tzu and eastern philosophies, and at the same time are strangely, inexplicably drawn back to the all-compelling reality of Jesus Christ.
\end{quotex}

\begin{quotex}
My teachings are very easy to understand and to practice; yet there is no one in the world who is able either to understand or to practice them. This is because my teachings have an originating principle and arise from an integrated system. This is not understood, so I am unknown.
\flright{\textsc{Lao Tzu}}
\end{quotex}

\begin{quotex}
Of the Logos, which is as I describe it, people always prove to be uncomprehending, both before and after they have heard of it. For although all things happen according to this Logos, people behave as if they have no experience, even when they experience such words and deeds as I explain, when I distinguish each thing according to its constitution and declare how it is. The rest of humanity fails to notice what they do after they wake up just as they forget what they do when asleep.
\flright{\textsc{Heraclitus}}
\end{quotex}



\flrightit{Posted on 2009-10-15 by Cologero}

\begin{center}* * *\end{center}

\begin{footnotesize}
\begin{sffamily}
\texttt{dagezhu on 2010-02-23 at 08:49 said:}

\url{https://orientem.blogspot.com/2009/01/in-beginning-was-tao.html}


has a similar argument


\hfill

\end{sffamily}
\end{footnotesize}

